\chapter*{Zusammenfassung}
\addcontentsline{toc}{chapter}{Zusammenfassung}
\markboth{Zusammenfassung}{Zusammenfassung}

\begin{otherlanguage}{ngerman}
Agentenbasierte Simulationen wie MATSim sind das Standardwerkzeug f\"ur die Bewertung st\"adtischer Verkehrspolitik, ben\"otigen jedoch Stunden pro Paris-Szenario; Graph-Neural-Network-Surrogate (GNN) reduzieren diese Laufzeit auf Sekunden.\ Genauigkeit allein gen\"ugt jedoch nicht f\"ur Entscheidungsunterst\"utzung: Vorhersagen k\"onnen sich au\ss{}erhalb der Trainingsverteilung unbemerkt verschlechtern.\ Diese Arbeit untersucht Unsicherheitsquantifizierung (UQ) f\"ur ein PointNetTransfGAT-GNN-Verkehrssurrogat am Pariser Stra\ss{}ennetz auf Basis einer festen Teilmenge von $1\,000$ aus $10\,000$ MATSim-Szenarien, die in allen elf Versuchen identisch verwendet wird.

Versuch~8 ($R^2 = 0{,}5957$ auf $3\,163\,500$ Testknoten) ist das prim\"are UQ-Basismodell.\ Monte-Carlo-Dropout liefert Spearman $\rho = 0{,}4820$ zwischen $\sigma$ und absolutem Fehler, doch die rohen Intervalle sind unkalibriert ($k_{95} = 11{,}66$ vs Gau\ss\ $1{,}96$); eine Regressions-$\sigma$-Skalierung mit $T = 2{,}887$ reduziert den Kuleshov-ECE um $90{,}5\,\%$ und bringt die $1\sigma$-Abdeckung innerhalb von $0{,}3$ Prozentpunkten an das Gau\ss-Ziel.\ Split-Konformalpr\"adiktion erreicht $\mathrm{PICP}_{90} = 90{,}02\,\%$ und $\mathrm{PICP}_{95} = 95{,}01\,\%$; adaptive Konformalpr\"adiktion verengt die bedingte Abdeckung von $[59{,}0\,\%,\,98{,}1\,\%]$ auf $[83{,}7\,\%,\,96{,}4\,\%]$ \"uber die $\sigma$-Dezile.\ Das Behalten der $50\,\%$ zuverl\"assigsten Vorhersagen reduziert den MAE um $41{,}2\,\%$ auf $2{,}32$ Fz/h.

Drei unsicherheitsbewusste Trainingsvarianten isolieren einen Designaspekt: Versuch~10 (CQR-Kopf, trainierbarer Backbone) bricht auf $R^2 = 0{,}4057$ ein, w\"ahrend Versuch~11 (gleicher Kopf und gleiche Lernrate, eingefrorener Backbone) alle sechs Kriterien bei $R^2 = 0{,}5835$ besteht -- der Vergleich st\"utzt das Einfrieren des MSE-trainierten Backbones innerhalb des untersuchten Rahmens und nicht als allgemeines Prinzip.\ Ein Deep Ensemble aus f\"unf Modellen erreicht die h\"ochste GNN-basierte, UQ-f\"ahige Punktgenauigkeit ($R^2 = 0{,}6841$, $+14{,}8\,\%$ \"uber T8) bei schw\"acherem UQ-Ranking ($\rho = 0{,}3997$).\ Nicht-GNN-Vergleichsmodelle (XGBoost $R^2 = 0{,}7414$, Random Forest $0{,}6612$) \"ubertreffen T8 in der Punktgenauigkeit.\ Der Beitrag dieser Arbeit ist orthogonal zur Punktgenauigkeit: Die hier evaluierte vorhersagespezifische UQ-Pipeline wird von den Baum-Vergleichsmodellen in ihrer Standardform nicht bereitgestellt; baumbasierte UQ-Erweiterungen wie Quantile Regression Forests und NGBoost sind plausible Konkurrenten und verbleiben als zuk\"unftige Arbeit.
\end{otherlanguage}
